%------------------------------------------------------------------------------%
\chapter{Prefácio}

% %------------------------------------------------------------------------------%
% \section*{Sobre a Disciplina}
% 
% Integração e Séries
% 
% Carga horária: 60 h.a.
% 
% Ementa: Integrais definidas: conceito, Teorema Fundamental do Cálculo e
% aplicações. Integrais indefinidas: conceito e métodos de integração. Integrais
% impróprias. Sequências e séries numéricas. Séries de potências, séries de
% Taylor e aplicações.
% 
% 
% Bibliografia Básica
% 1. THOMAS, George B. Cálculo. 11. ed. São Paulo: Pearson, 2008. v. 1 e v. 2.
% 2. STEWART, James. Cálculo. 5. ed. São Paulo: Thomson, 2003. v. 1 e v. 2.
% 3. FLEMMING, Diva Marília; GONÇALVES, Mirian Buss. Cálculo A: funções, limite, derivação e integração. São Paulo: Prentice-Hall, 2007.
% 
% Bibliografia Complementar
% 1. EDWARDS JR., C. H.; PENNEY, David E. Cálculo com geometria analítica. Rio de Janeiro: Prentice-Hall, 1994. v. 1 e v.2.
% 2. SWOKOWSKI, Earl W. Cálculo com geometria analítica. 2. ed. São Paulo: Makron Books, 1995. v.1 e v.2.
% 3. SIMMONS, George F. Cálculo com geometria analítica. São Paulo: Pearson Makron Books, 1988. v. 1 e v.2.
% 4. LEITHOLD, Louis. O cálculo com geometria analítica. 3. ed. São Paulo: Harbra, 1994. v. 1 e v.2.
% 5. BOULOS, P. Cálculo diferencial e integral. São Paulo: Makron Books,1999. v. 1.

\vspace{3cm}
\begin{observation}
  Este texto ainda está em construção e, portanto, está incompleto e contém erros,
  ele não é um substituto para as aulas e livros da bibliografia.
\end{observation}
\vspace{3cm}

Essa apostila foi idealizada para apresentar o conteúdo programático previsto para
a disciplina de \emph{Integrais e Séries} no \emph{CEFET-MG},
cuja ementa contém os tópicos:
\begin{enumerate}[nosep, topsep=-0.5\baselineskip]
  \item integrais definidas;
  \item integrais Indefinidas;
  \item séries numéricas e
  \item séries de potências.
\end{enumerate}

O conteúdo dessa apostila está distribuído nos seguintes capítulos:
\vspace{-0.5\baselineskip}
\begin{enumerate}[nosep]
  \item \nameref{cap:funcoes_derivadas}
  \item \nameref{cap:integracao}
  \item \nameref{cap:areas_volumes}
  \item \nameref{cap:tecnicas_integracao}
  \item \nameref{cap:sequencias_numericas}
  \item \nameref{cap:series_numericas}
  \item \nameref{cap:series_potencias}
  \item \nameref{cap:series_de_taylor}
  \item \nameref{cap:usos_taylor}
\end{enumerate}
Cada capítulo apresenta a teoria, exemplos e exercícios sobre o tópico correspondente.
Sempre que possível apresentamos também exemplos e ilustrações implementadas em
\href{https://www.python.org}{Python} disponibilizadas como
\textit{notebooks} no
\href{https://colab.research.google.com}{Google Colaboratory} ou Colab.
Nenhuma das atividades em Python é necessária para o estudo da apostila,
elas foram criadas apenas como atividades complementares para os alunos interessados.
A seguir está o \textit{link} para o \textit{notebook} que apresenta as bibliotecas
\href{https://numpy.org}{NumPy}, usada para computação numérica, e
\href{https://www.sympy.org}{SymPy}, usada para computação simbólica.

\OnlineBox{https://colab.research.google.com/drive/1XV6A0l2plJHVPVfb2xBVDmJA4Hkm3MhL?usp=share_link}
{
  Notebook introduz as bibliotecas NumPy e SymPy.
}

Como os textos matemáticos costumam ser muito densos em informação,
não absorvemos todas as sutilezas apenas em uma leitura. Por isso, fazer exercícios é
parte fundamental do processo de aprendizado, ao manipularmos os conceitos no processo
de resolução do exercício somos forçados a utilizar os resultados e a explorar suas
sutilezas. A regra geral é que cada passagem da resolução precisa ser embasada em um
resultado verdadeiro explicitamente descrito antes, essa é a razão pela qual os textos
numeraram os resultados e fórmulas. Mesmo que não escrevamos essas referências,
precisamos estar plenamente conscientes delas.

Uma sugestão é utilizar os exemplos contidos no texto como exercícios, tente resolver
o problema proposto antes de ler a resolução apresentada. Depois dessa tentativa leia
a resolução e a compare criticamente à sua resposta.
Seu resultado coincide com o apresentado? Se sim verifique se você apresentou todos os
argumentos necessários para embasar sua resposta.
Quando os resultados ou argumentos não coincidirem revise
sua resolução e verifique se ela corresponde a uma solução alternativa ou contém
algum erro. Nas disciplinas de Matemática, na maioria das vezes, a argumentação
é muito mais importante do que o resultado.

Nesse texto foi inserida uma lista de exercícios após cada seção e uma lista de revisão
no final de cada capítulo. Ao estudar uma seção, teste sua compreensão resolvendo alguns
exercícios da seção. Ao fazer uma revisão ou se preparar para uma avaliação, resolva
os exercícios da revisão. O fato deles não estarem organizados por temas introduz o
processo de identificação do problema e escolha da teoria adequada para resolvê-los.

Alguns exercícios possuem resposta no final da apostila, nesses casos o
\textit{link} \textcolor{cpLinks}{[resp]} no enunciado do exercício
leva para sua resposta. Para retornar ao enunciado clique no número da resposta.
Em alguns casos a resposta consiste apenas do resultado final, em outros toda a
resolução está escrita.

Alguns resultados matemáticos necessários para a disciplina, como números complexos e
funções hiperbólicas, são apresentados sucintamente no apêndice
\nameref{cap:Conteudo_Complementar}.
O apêndice \nameref{cap:Recursos_Online} apresenta diversos materiais complementares,
como videoaulas e atividades \textit{online}, que podem ser acessados gratuitamente
e contribuir para o estudo de cada tópico.
O \nameref{cap:indice_remissivo} traz \textit{links} para os principais termos
discutidos dentro do texto.

%------------------------------------------------------------------------------%
